\subsection{整中心滤过}

回顾 (no. 1， 注) 群 G 上的滤过 \((\mathrm{G}_{\alpha})\) 称为整的，如果 \(\mathrm{G}_{\alpha} = \mathrm{G}_{n}\) （对每个整数 \(n\) 以及一切 \(\alpha \in ]n - 1,n]\) ）。给定群 G 上的一个整中心滤过等价于给定 G 的子群的一个序列 \((\mathrm{G}_n)_{n\geqslant 1}\) ，它满足条件

\begin{enumerate}
\def\labelenumi{(\roman{enumi})}
\tightlist
\item
\end{enumerate}

\[
\mathrm{G} _ {1} = \mathrm{G}\tag{i}
\]

\[
\mathrm{G} _ {n} \supset \mathrm{G} _ {n + 1} \quad \text { for   all } n \geqslant 1\tag{ii}
\]

\[
\left(\mathrm{G} _ {m}, \mathrm{G} _ {n}\right) \subset \mathrm{G} _ {m + n} \quad \text { for } m \geqslant 1 \text { and } n \geqslant 1.\tag{iii}
\]

对每个整数 \(n \geqslant 1\) ， \(G_n\) 是 \(G\) 的一个正规子群，并且商 \(\mathrm{gr}_n(G) = G_n / G_{n+1}\) 是交换的。取商后，映射 \((x,y) \mapsto (x,y) = x^{-1}y^{-1}xy\) （从 \(G_m \times G_n\) 到 \(G_{m+n}\) 中）使我们能在 \(\mathrm{gr}(G) = \bigoplus_{n \geqslant 1} \mathrm{gr}_n(G)\) 上定义一个 \(N\) 型分次李代数结构（在环 \(Z\) 上）。

回顾 (《代数》， 第 I 章， § 6， no. 3， 定义 5)，群 \(\mathbf{G}\) 的下中心列由

\[
\mathrm{C} ^ {1} \mathrm{G} = \mathrm{G}, \quad \mathrm{C} ^ {n + 1} = (\mathrm{G}, \mathrm{C} ^ {n} \mathrm{G}) \quad \text { for } n \geqslant 1.\tag{28}
\]

定义。相应的滤过称为 G 的下中心滤过。

命题 4. (i) G 的下中心列是 G 上的一个整中心滤过。

\begin{enumerate}
\def\labelenumi{(\roman{enumi})}
\setcounter{enumi}{1}
\tightlist
\item
  若 \((\mathrm{G}_n)_{n\in \mathbf{N}^*}\) 是 \(\mathbf{G}\) 上的一个整中心滤过，则 \(\mathbf{C}^n\mathbf{G}\subset \mathbf{G}_n\) （对一切 \(n\in \mathbf{N}^*\) ）。
\end{enumerate}

断言 (i) 已在《代数》第 I 章 § 6 no. 3 公式 (7) 中证明。

我们对 \(n\) 作归纳证明 (ii)； \(\mathrm{C}^1\mathrm{G} = \mathrm{G} = \mathrm{G}_1\) ；对 \(n > 1\) ，

\[
\mathrm{C} ^ {n} \mathrm{G} = (\mathrm{G}, \mathrm{C} ^ {n - 1} \mathrm{G}) \subset (\mathrm{G}, \mathrm{G} _ {n - 1}) \subset \mathrm{G} _ {n}.
\]

命题 5. 设 \(\mathbf{G}\) 是一个群， \(\operatorname{gr}(\mathbf{G})\) 是分次李 \(\mathbf{Z}\) -代数，它与 \(\mathbf{G}\) 上的下中心滤过相伴。那么 \(\operatorname{gr}(\mathbf{G})\) 由 \(\operatorname{gr}_1(\mathbf{G}) = \mathbf{G}/(\mathbf{G},\mathbf{G})\) 生成。

设 L 是 \(\mathrm{gr}(\mathrm{G})\) 的由 \(\mathrm{gr}_{1}(\mathrm{G})\) 生成的李子代数；我们对 n 作归纳证明 \(\mathrm{L} \supset \mathrm{gr}_{n}(\mathrm{G})\) ，断言对 n = 1 是平凡的。设 n > 1 且 \(\mathrm{L} \supset \mathrm{gr}_{n-1}(\mathrm{G})\) 。由于 \(\mathrm{C}^{n}\mathrm{G} = (\mathrm{G}, \mathrm{C}^{n-1}\mathrm{G})\) ， \(\mathrm{gr}(\mathrm{G})\) 上李代数法则的构造立即表明

\[
\operatorname{gr} _ {n} (\mathrm{G}) = \left[ \operatorname{gr} _ {1} (\mathrm{G}), \operatorname{gr} _ {n - 1} (\mathrm{G}) \right] \subset \mathrm{L}.
\]

上面的证明表明，李代数 \(\mathbf{gr}(\mathbf{G})\) 的下中心列 (§ 2， no. 7) 由下式给出：

\[
\mathcal {C} ^ {n} (\mathrm{gr} (G)) = \sum_ {m \geqslant n} \mathrm{gr} _ {m} (G).\tag{29}
\]

注. 设 \(k\) 是一个环， \(n\) 是一个整数 \(>0\) ，A 是由 \(n\) 行 \(n\) 列、元素在 \(k\) 中的下三角矩阵组成的集合。对 \(p \geqslant 0\) ，令 \(\mathbf{A}_p\) 为 \((x_{ij}) \in \mathbb{A}\) 组成的集合，其中 \(x_{ij} = 0\) （对 \(i - j < p\) ）。那么 \(\mathrm{A}_0 = \mathrm{A}\) 且 \(\mathrm{A}_p\mathrm{A}_q \subset \mathrm{A}_{p + q}\) 。设 \(\Gamma_p = 1 + \mathrm{A}_p\) 。则 \(\Gamma_1\) 是 \(\mathbf{GL}(n,k)\) 的一个子群，称为 \(n\) 阶的严格下三角群（在 \(k\) 上）。由 no. 5 的命题 2， \((\Gamma_p)\) 是 \(\Gamma_1\) 上的一个整滤过。由于 \(\Gamma_n = \{1\}\) ，我们看出群 \(\Gamma_1\) 是幂零的 (《代数》， 第 I 章， § 6， no. 3， 定义 6)。

